\documentclass[12pt]{article}
\usepackage[utf8]{inputenc}
\usepackage[T1]{fontenc}
\usepackage[margin=1in]{geometry}
\usepackage{setspace}
\usepackage{xcolor}
\usepackage[colorlinks=true, urlcolor=blue!60!black, linkcolor=blue!60!black]{hyperref}
\setlength{\parindent}{2em}
\setlength{\parskip}{1.25em}
\renewcommand{\baselinestretch}{1.0}
\begin{document}
\begin{center}
{\LARGE \textit{Reynolds v. United States}, 98 U.S. 145 (1879)\par}
\vspace{0.8em}
{\large Case Brief --- Four-Page Edition\par}
\end{center}
\vspace{1em}

\section{Memory Jogger}

In its first major interpretation of the Free Exercise Clause, the Supreme Court drew a sharp line between religious belief and religious conduct, holding that sincere religious duty does not excuse violating a generally applicable federal criminal prohibition on polygamy (\href{https://supreme.justia.com/cases/federal/us/98/145/}{Justia}).

\section{Detailed Case Facts}

George Reynolds, a Utah Territory resident and devout Latter-day Saint, married Amelia Jane Schofield while his first wife, Mary Ann Tuddenham, was living, in a ceremony authorized by Church officials. He proved that his Church taught a divinely revealed duty for men to practice polygamy, on pain of ``damnation in the life to come.'' The Morrill Anti-Bigamy Act of 1862 (Rev. Stat. \S~5352) made marrying another while having a living spouse in a federal territory a crime punishable by a \$500 fine and five years in prison. The prosecution's witness Schofield could not be served; Reynolds told the deputy marshal ``that will be for you to find out,'' and her testimony from an earlier trial was admitted over a Confrontation Clause objection. The judge refused an instruction to acquit if Reynolds acted from religious duty and urged jurors to consider the ``innocent victims of this delusion.''

\section{Procedural History}

An earlier indictment was held invalid. Reynolds was reindicted by a fifteen-person grand jury in the Third Judicial District of Utah Territory, convicted, and sentenced to two years at hard labor and a \$500 fine. The Supreme Court of the Territory affirmed. On writ of error the U.S. Supreme Court affirmed in January 1879; on rehearing it reversed only to strike ``at hard labor,'' which the statute did not authorize.

\section{Judicial Votes}

Affirmed (sentence later corrected), \textbf{9--0} on free exercise; \textbf{8--1} on the Confrontation Clause issue.
\begin{itemize}
\item \textbf{Majority:} Chief Justice Waite, joined by Justices Clifford, Swayne, Miller, Strong, Bradley, Hunt, and Harlan.
\item \textbf{Concurring in part and dissenting in part:} Justice Field (dissenting only on Schofield's former testimony).
\end{itemize}

\section{Holding}

Congress has no legislative power over mere religious belief, but may punish overt acts that violate social duties or subvert good order; religious belief or duty is no defense to a generally applicable bigamy statute and does not negate criminal intent. The Court also held that territorial law governed grand jury size, that juror rulings stand absent manifest error, that practicing polygamists could be excluded from the jury, and that a defendant who wrongfully procures a witness's absence forfeits his confrontation objection. Waite's reliance on Jefferson's ``wall of separation'' and his remarks about polygamy among ``Asiatic and of African people'' were influential dicta.

\section{Analysis of Opinions}

\textbf{Majority (Waite).} Because the Constitution does not define ``religion,'' Waite turned to Virginia history: Madison's \textit{Memorial and Remonstrance}, Jefferson's statute allowing interference ``when principles break out into overt acts against peace and good order,'' and Jefferson's Danbury letter saying legislative powers ``reach actions only, and not opinions.'' Treating that letter as authoritative, he announced the belief-action line. Polygamy had been a crime in England and the colonies, and Virginia reenacted its ban in 1788. Marriage is a civil contract on which society is built, and, citing Francis Lieber, polygamy fosters ``the patriarchal principle'' and ``stationary despotism.'' Religiously mandated human sacrifice or widow-burning could not be excused either. Because Reynolds knowingly married again, intent was implied. On confrontation, drawing on \textit{Lord Morley's Case}, no one may profit from his own wrong in procuring a witness's absence. Territorial courts are not Article III ``circuit and district courts,'' so Utah's fifteen-person grand jury statute governed; juror Charles Read's unexpressed impression showed no manifest partiality; and the judge's remarks about polygamy's ``innocent victims'' were a permissible reminder of the jury's duty rather than an appeal to prejudice.

\textbf{Field (in part).} Joined the free exercise holding but found no adequate foundation that Reynolds procured Schofield's absence; the cited authorities ``establish conclusively the exact reverse.''

\section{Examples: Future Applications}

\begin{enumerate}
\item \textbf{Same-Side: Sacramental Drug Use.} Members of a sect ingest an illegal hallucinogen as a sacrament. The State may punish the act though not the belief, as \textit{Employment Division v. Smith} (1990) held.
\item \textbf{Same-Side: Social Security Taxes.} An employer whose faith forbids government insurance refuses payroll taxes. Religious duty does not excuse compliance with a neutral tax (\textit{United States v. Lee}).
\item \textbf{Opposite-Side: Punishing Belief.} A territorial law punishes professing or preaching plural marriage. Congress has no power over ``mere religious belief and opinions.''
\item \textbf{Opposite-Side: Targeting a Minority Ritual.} A city bans animal sacrifice by one faith while exempting secular slaughter. Even \textit{Reynolds} and \textit{Smith} forbid targeting religion (\textit{Lukumi}).
\item \textbf{Fence-Sitter: Amish Schooling.} A neutral compulsory-attendance law is applied to a separatist agrarian community that ends schooling after eighth grade. Strict \textit{Reynolds} favors the State, but \textit{Wisconsin v. Yoder} (1972) granted an exemption.
\end{enumerate}

\section{Critique}

\textbf{Originalist and textualist critiques.} Michael McConnell faults Waite for treating Jefferson, who was in France during the drafting, as the authoritative voice while overlooking Madison and state provisos protecting exercise unless it disturbed peace or safety. Textualists note that ``free \textit{exercise}'' denotes practice, not mere belief. Clark Lombardi reads Waite as asking whether polygamy fell within traditional limits of peace and good order.

\textbf{Progressive, doctrinal, and historical critiques.} Critics see Protestant majoritarianism and anti-Mormon hostility in Waite's culturally biased claims. The Court moved to strict scrutiny in \textit{Sherbert} and \textit{Yoder} before Justice Scalia revived the ``law unto himself'' warning in \textit{Smith}. Defenders of \textit{Reynolds} and \textit{Smith} argue that exemptions threaten civil rights and third parties. The case pits religious conscience against equal rule of law.

\section{Key Quotations}

\begin{quote}
``Can a man excuse his practices to the contrary because of his religious belief? To permit this would be to make the professed doctrines of religious belief superior to the law of the land, and, in effect, to permit every citizen to become a law unto himself. Government could exist only in name under such circumstances.''
\end{quote}
\noindent --- Chief Justice Waite, 98 U.S. at 166--167.

\begin{quote}
``Laws are made for the government of actions, and while they cannot interfere with mere religious belief and opinions, they may with practices.''
\end{quote}
\noindent --- Chief Justice Waite, 98 U.S. at 166.

\begin{quote}
``Congress was deprived of all legislative power over mere opinion, but was left free to reach actions which were in violation of social duties or subversive of good order.''
\end{quote}
\noindent --- Chief Justice Waite, 98 U.S. at 164.

\end{document}
